%--------------------------------------------------------------------------------
\pagebreak
\unnumberedSection{BR - Branch Register}

%--------------------------------------------------------------------------------
\begin{iPageDescEntry}{Syntax}
\texttt{BR RegB [, RegR ]}
\end{iPageDescEntry}

%--------------------------------------------------------------------------------
\begin{iPageDescEntry}{Format}
The \texttt{BR} instruction uses the following instruction format: \\

\begin{flushleft}
\drawInstrFormat
    {0,1,3,5,6.5,8.5,16}
    {\OpCodeGrpBR, \OpCodeBR, RegR, 0, RegB, 0}
\end{flushleft}
\end{iPageDescEntry}

%--------------------------------------------------------------------------------
\begin{iPageDescEntry}{Description}

The \texttt{BR} instruction performs an unconditional branch to a target address computed relative to the current instruction address (\texttt{IA}). The target address is formed by taking the contents of \texttt{RegB}, left shifting the value by two, and adding the result to the instruction address offset portion.

Optionally, a return link may be written to \texttt{RegR}. When specified, \texttt{RegR} receives the address of the next sequential instruction, allowing the \texttt{BR} instruction to be used for procedure calls or control-flow transfers that require a return address.

\end{iPageDescEntry}

%--------------------------------------------------------------------------------
\begin{iPageDescOpEntry}{Operation}
\begin{lstlisting}[style=iPageOpStyle]
if ( instr.[RegR] != 0 ) {

    RegR.[51:32] <- PSR.[IA-REG];
    RegR.[31:0]  <- addOfs32( PSR.[IA-OFS], 4 );
}

PSR.[IA-OFS] <- addOfs32( PSR.[IA-OFS], RegB << 2 ));
\end{lstlisting}
\end{iPageDescOpEntry}

%--------------------------------------------------------------------------------
\begin{iPageDescEntry}{Exceptions}
Taken Branch Trap \\
Instruction TLB Miss Trap
\end{iPageDescEntry}

%--------------------------------------------------------------------------------
\begin{iPageDescEntry}{Notes}
None
\end{iPageDescEntry}